
\section{Cross-References}\fbeg{}\begin{frame}{\ft{What's \q{Cross-Reference Logic}?}}


\begin{ftxt}[.93\textwidth]

Note that for a PDF rendering library, each PDF page is just a picture, 
which can be manipulated and overlaid with generic graphics 
functionality.  The window coordinates of textual elements 
can be obtained from PDF annotations or by pairing PDF files 
with metadata compiled using techniques such as \pdfsp{13pt}.  
This means that context menus activated over a PDF page can 
have different options depending on the contents of 
words, sentences, or annotations present near or around 
the mouse-event position.  Also, the PDF display
can be embedded in GUI windows with additional buttons and 
controls.  Whether via context menus, toolbars, or other 
controls, Research Object Applications can link PDF 
windows programmatically with other windows showing videos, 
digital maps, computer simulations, 3D graphics, and so forth.
PDF libraries such as Poppler or XPDF can be included 
directly in Research Objects, so PDF viewers   
can be customized for individual data set and 
shared directly as source-code files. 


\end{ftxt}


\end{frame}

